\begin{theorem}
Следующие утверждения эквивалентны:

\begin{enumerate}
    \item $\sk{k} = \sk{k + 1}$
    \item $\sk{k} = \pk{k}$
    \item $\ph = \sk{k}$ (полиномиальная иерархия схлопывается после $k$-го уровня)
\end{enumerate}
\end{theorem}

\begin{proof}
Из последнего утверждения первые два следуют по определению (второе из замкнутости $PH$ относительно дополнения), поэтому интересен вывод в обратном направлении. Вначале докажем эквивалентность первых двух. Затем по индукции докажем, что $\sk{k} = \sk{k+1}$. Отсюда легко следует PH = $\sk{k}$.
Действительно, пусть $\sk{k} = \pk{k}$, а $A \in \sk{k + 1}$. Докажем, что $A \in \sk{k}$. Для определённости, пусть $k$ нечётно. Значит, 
$x \in A \lra \exists y_1 \forall y_2 \exists y_3 \ldots \forall y_{k + 1} V(x, y_1, y_2, y_3, \ldots, y_{k + 1}) = 1$. Тогда множество $\{ (x, y_1) | \forall y_2 \exists y_3 \ldots \forall y_{k + 1} V(x, y_1, y_2, y_3, \ldots, y_{k + 1}) = 1\}$ лежит в $\pk{k}$. По предположению оно лежит и в $\sk{k}$,
т.е. принадлежность ему эквивалентна $\exists z_2 \forall z_3 \ldots \exists z_{k+1} W(x, y_1, z_2, \ldots, z_{k+1})=1$ для некоторого $W$. Таким образом, $x \in A \lra \exists y_1 \exists z_2 \forall z_3 \ldots \exists z_{k + 1} W(x, y_1, z_2, z_3, \ldots, z_{k + 1}) = 1$. А это уже формула из $\sk{k}$, поскольку кванторы по $y_1$ и $z_2$ можно объединить
в один.

Обратно, пусть $\sk{k} = \sk{k + 1}$. Тогда $\pk{k} \subset \sk{k + 1}$, следовательно, $\pk{k} \subset \sk{k}$. Отсюда верно и обратное включение: если $B \in \sk{k}$, то $\overline{B} \in \pk{k}$, откуда $\overline{B} \in \sk{k}$, откуда $B \in \pk{k}$. Значит, $\sk{k} = \pk{k}$.

Далее рассуждение проходит похожим образом: если $B \in \pk{k + 1}$, то $\overline{B} \in \sk{k + 1}$, откуда по доказанному $\overline{B} \in \sk{k}$, откуда $B \in \pk{k}$. Отсюда $\sk{k + 1} = \pk{k + 1} = \sk{k}$, откуда по индукции $\sk{k + l} = \pk{k + l} = \sk{k}$ и, как следствие, $\ph = \sk{k}$.
\end{proof}

\begin{statement}
Если в классе \ph существует полная задача относительно сводимости по Карпу, то \ph = \sk{k} для некоторого $k$.
\end{statement}

\begin{proof}
Действительно, если язык $A$ является \ph-полным, то он, в частности, лежит в \ph и потому лежит в \sk{k} для некоторого $k$. Но раз любой другой $B \in \ph$ сводится к $A$, то $B$ также лежит в \sk{k}, откуда \ph = \sk{k}, что и требовалось.
\end{proof}
